\subsection*{CRAAP-Test zur Quellenüberprüfung}

\subsubsection*{CRAAP-Test Quelle \cite{zahorulko_ai_2025}}
\begin{longtable}{|>{\raggedright\arraybackslash}p{3cm}|>{\raggedright\arraybackslash}p{6cm}|>{\raggedright\arraybackslash}p{7cm}|}

\caption{CRAAP-Test Quelle \cite{zahorulko_ai_2025}} \label{tab:craap-zahorulko-ai} \\ 
\hline
\textbf{Kriterium} & \textbf{Fragestellungen} & \textbf{Bewertung / Notizen} \\
\hline
\endfirsthead

\hline
\textbf{Kriterium} & \textbf{Fragestellungen} & \textbf{Bewertung / Notizen} \\
\hline
\endhead

\textbf{C – Currency (Aktualität)}  &
- Wann wurde die Quelle veröffentlicht oder zuletzt aktualisiert?  \newline
- Ist die Information noch aktuell genug? [cite: 16] \newline
- Gibt es neuere Erkenntnisse? [cite: 16] \newline
- Funktionieren die referenzierten Links noch? [cite: 16] &
% Hier deine Notizen eintragen
\\ \hline

\textbf{R – Relevance (Relevanz)} [cite: 17] &
- Passt die Quelle zu deinem Thema oder deiner Fragestellung? [cite: 17] \newline
- Ist die Information auf dein Niveau abgestimmt? [cite: 18] \newline
- Ist die Quelle primär oder sekundär? [cite: 18] &
% Hier deine Notizen eintragen
\\ \hline

\textbf{A – Authority (Autorität)} [cite: 19] &
- Wer ist der Autor / die Institution? [cite: 19] \newline
- Welche Qualifikationen oder Fachkenntnisse sind vorhanden? [cite: 19] \newline
- Ist die Quelle vertrauenswürdig (z.B. Universität, Fachjournal)? [cite: 20] &
% Hier deine Notizen eintragen
\\ \hline

\textbf{A – Accuracy (Genauigkeit)} [cite: 21] &
- Ist die Information belegt durch Quellen oder Beweise? [cite: 21] \newline
- Gibt es Hinweise auf Fehler oder Voreingenommenheit? [cite: 21] \newline
- Wurde die Quelle gegengelesen (Peer Review)? [cite: 22] &
% Hier deine Notizen eintragen
\\ \hline

\textbf{P – Purpose (Zweck)} [cite: 23] &
- Was ist das Ziel der Quelle (informieren, verkaufen, überzeugen)? [cite: 23] \newline
- Gibt es Anzeichen für Meinung oder Werbung? [cite: 23] \newline
- Ist die Sprache sachlich oder emotional gefärbt? [cite: 24] &
% Hier deine Notizen eintragen
\\ \hline

\end{longtable}
\newpage